\section{Ejemplos por materia}\label{sec:ejemplos}

El mismo inventario sirve para cualquier curso. En cada ejemplo aparecen primero
el apunte y después su lectura en español. Estos ejemplos usan sobre todo
conectores, para que se lean con facilidad; el estilo completo, con raíces y
estilo telegráfico, está en las diez clases de \path{compilado/apuntes.pdf}.

\subsection{Matemática}
\begin{apunte}[Cálculo · límites]
\linea{}{\ks{K09} \(f\) continua en \(a\): \(\lim_{x\to a}f(x)=f(a)\)}
\linea{\ks{K31}}{derivable \(\Rightarrow\) continua}
\linea{}{\(\neg\)[continua \(\Rightarrow\) derivable] \quad \ks{K26} \(|x|\) en 0}
\end{apunte}
\lectura{Una función es continua en \(a\) si su límite en \(a\) coincide con su
valor. Importante: si es derivable, es continua. Lo contrario no se cumple; por
ejemplo, \(|x|\) en cero.}

\subsection{Computación}
\begin{apunte}[Teoría de la computación · jerarquía de Chomsky]
\linea{}{regulares \(\subset\) libres de contexto \(\subset\) decidibles}
\linea{}{AFD \(\overset{\text{\tiny acepta}}{\longrightarrow}\) regulares \quad autómata de pila \(\overset{\text{\tiny acepta}}{\longrightarrow}\) libres de contexto}
\linea{\ks{K31}}{\ks{K16} \ks{K30} Turing \ks{K07} 1936: problema de la parada \(\notin\) decidibles}
\linea{\ks{K32}}{¿por eso ningún compilador detecta todos los bucles infinitos?}
\end{apunte}
\lectura{Todo lenguaje regular es libre de contexto, y todo libre de contexto es
decidible. Los AFD aceptan los regulares; los autómatas de pila, los libres
de contexto. Importante: según Turing (1936), el problema de la parada
no es decidible. Idea mía: por eso ningún compilador puede detectar todos los
bucles infinitos.}

\subsection{Física}
\begin{apunte}[Dinámica]
\linea{}{\ks{K16} Newton: \(F=ma\)}
\linea{}{\(\uparrow F\) \(\wedge\) masa fija \(\Rightarrow\) \(\uparrow a\)}
\linea{}{fricción \ks{K23} calor \quad energía cinética \ks{K24} térmica}
\linea{\ks{K34}}{ejercicios 4--9, cap.~3}
\end{apunte}
\lectura{Según Newton, la fuerza es masa por aceleración: si la fuerza aumenta
y la masa no cambia, la aceleración aumenta. La fricción causa calor: la
energía cinética se convierte en térmica. Pendiente: ejercicios 4 al 9.}

\subsection{Química}
\begin{apunte}[Reacciones]
\linea{}{\(\mathrm{CH_4 + 2\,O_2}\) \ks{K24} \(\mathrm{CO_2 + 2\,H_2O}\)}
\linea{}{combustión \(\subset\) reacción redox}
\linea{}{ácido \ks{K29} base \quad pH \(<7\) \ks{K25} pH \(>7\)}
\end{apunte}
\lectura{El metano con oxígeno se convierte en dióxido de carbono y agua. La
combustión es un tipo de reacción redox. Ácido y base son extremos opuestos:
pH menor que 7, en cambio mayor que 7.}

\subsection{Biología}
\begin{apunte}[Célula]
\linea{}{mitocondria \ks{K27} célula eucariota}
\linea{}{\ks{K04} respiracn celular: glucosa \ks{K24} ATP}
\linea{\ks{K33}}{\ks{K16} Margulis: mitocondria \(\approx\) antigua bacteria}
\end{apunte}
\lectura{La mitocondria es parte de la célula eucariota. En la respiración
celular, la glucosa se convierte en ATP. Hipótesis, según Margulis: la
mitocondria viene de una antigua bacteria.}

\subsection{Geografía y ambiente}
\begin{apunte}[Clima]
\linea{}{\ks{K08} Perú: corriente de Humboldt \ks{K23} costa árida}
\linea{}{Niño \ks{K23} \(\uparrow\) temperatura del mar \ks{K23} lluvias en la costa}
\linea{}{clima \(\neq\) tiempo atmosférico}
\end{apunte}
\lectura{En el Perú, la corriente de Humboldt causa que la costa sea árida. El
Niño aumenta la temperatura del mar, lo que causa lluvias en la costa. El clima
no es lo mismo que el tiempo atmosférico.}

\subsection{Psicología}
\begin{apunte}[Desarrollo]
\linea{}{\ks{K16} \ks{K30} Piaget: sensoriomotora \ks{K21} preoperacional \ks{K21} concreta \ks{K21} formal}
\linea{}{conservacn \(\in\) logros de la etapa concreta}
\linea{\ks{K19}}{¿las etapas son iguales en todas las culturas?}
\end{apunte}
\lectura{Según Piaget, el desarrollo pasa por cuatro etapas en ese orden. La
conservación es uno de los logros de la etapa concreta. Duda: ¿las etapas son
iguales en todas las culturas?}

\subsection{Economía}
\begin{apunte}[Mercado]
\linea{}{\(\uparrow\) precio \ks{K23} \(\downarrow\) cantidad demandada}
\linea{}{\ks{K09} equilibrio: oferta \(=\) demanda}
\linea{\ks{K35}}{\(\uparrow\) demanda \(\wedge\) oferta fija \(\Rightarrow\) \(\uparrow\) precio}
\end{apunte}
\lectura{Si sube el precio, baja la cantidad demandada. El equilibrio es el punto
donde la oferta y la demanda son iguales. Resumen: si la
demanda sube y la oferta no cambia, el precio sube.}

\subsection{Historia}
\begin{apunte}[Independencia del Perú]
\linea{}{\ks{K07} 1821 \ks{K08} Lima \ks{K30} San Martín: proclamacn}
\linea{}{proclamacn \ks{K21} Ayacucho (\ks{K07} 1824) \ks{K23} fin del virreinato}
\linea{\ks{K32}}{1821 \(\neq\) independencia efectiva}
\end{apunte}
\lectura{En 1821, en Lima, San Martín proclamó la independencia. Después vino
la batalla de Ayacucho (1824), que causó el fin del virreinato. Idea mía: 1821
no es la fecha de la independencia efectiva.}

\subsection{Derecho y política}
\begin{apunte}[Estado]
\linea{}{poder ejecutivo, legislativo, judicial \ks{K27} Estado}
\linea{}{\ks{K16} Montesquieu: separación de poderes \ks{K23} \(\downarrow\) abuso}
\linea{}{ley \(\neq\) decreto: quién la dicta}
\end{apunte}
\lectura{Los poderes ejecutivo, legislativo y judicial son partes del Estado.
Según Montesquieu, separarlos reduce el abuso de poder. Una ley no es lo mismo
que un decreto: los dicta un órgano distinto.}

\subsection{Lengua y literatura}
\begin{apunte}[Figuras retóricas]
\linea{}{metáfora \(\subset\) figura retórica \quad \ks{K26} «tus ojos son luceros»}
\linea{}{metáfora \ks{K25} símil: el símil usa «como»}
\linea{}{\ks{K30} Vallejo \(\in\) vanguardia \quad \ks{K26} \emph{Trilce}, \ks{K07} 1922}
\end{apunte}
\lectura{La metáfora es un tipo de figura retórica; por ejemplo, «tus ojos son
luceros». Se diferencia del símil en que este usa «como». Vallejo es uno de los
poetas de la vanguardia; por ejemplo, \emph{Trilce}, de 1922.}

\subsection{Arte y música}
\begin{apunte}[Armonía]
\linea{}{acorde mayor \ks{K09} fundamental \(+\) 3.ª mayor \(+\) 5.ª justa}
\linea{}{mayor \ks{K25} menor: la tercera}
\linea{}{tensión \ks{K21} resolucn \quad \ks{K26} V \ks{K21} I}
\end{apunte}
\lectura{Un acorde mayor es la fundamental con su tercera mayor y su quinta
justa. Se diferencia del menor en la tercera. La tensión va seguida de la
resolución; por ejemplo, el acorde V seguido del I.}
